% ══════════════════════════════════════════════════════════════════════════════
% Tables
% ══════════════════════════════════════════════════════════════════════════════

\begin{table}[p]
\centering
\caption{\textbf{Demographics and clinical characteristics of the 23 patients
with complete longitudinal data across all three visits.} Values are reported as
mean~$\pm$~SD, median (range), or $n$~(\%) as appropriate.}
\label{tab:demographics}
\begin{tabular}{lccc}
\toprule
Characteristic & Mean $\pm$ SD & Median (Range) & $n$ (\%) \\
\midrule
Age at enrollment (yrs) & 30.2 $\pm$ 5.1 & 29 (20--41) & \\
\addlinespace
Race & & & \\
\quad White & & & 22 (96) \\
\quad Not disclosed & & & 1 (4) \\
\addlinespace
Prepregnancy BMI & 26.6 $\pm$ 5.3 & 25 (19.0--37.8) & \\
\addlinespace
Parity & & 0 (0--2) & \\
\quad Nulliparous & & & 16 (70) \\
\quad Parous & & & 7 (30) \\
\addlinespace
Enrollment condition & & & \\
\quad Healthy nulliparous & & & 14 (61) \\
\quad PCOS & & & 3 (13) \\
\quad Prior PE & & & 6 (26) \\
\addlinespace
Cycle day at V1 & 9.2 $\pm$ 3.6 & 10 (3--14) & \\
\addlinespace
Gestational age at V2 (days) & 89.0 $\pm$ 5.5 & & \\
Gestational age at V3 (days) & 217.0 $\pm$ 4.0 & & \\
\midrule
\multicolumn{4}{l}{\textit{Enrollment cohort $\times$ pregnancy outcome}} \\
\midrule
 & Uncomplicated & Developed PE & Total \\
\midrule
Healthy nulliparous & 12 & 2 & 14 \\
Prior PE & 4 & 2 & 6 \\
PCOS & 2 & 1 & 3 \\
\midrule
Total & 18 & 5 & 23 \\
\bottomrule
\end{tabular}
\end{table}

\begin{table}[p]
\centering
\caption{Per-visit mean relative error (MRE) between real and SA-generated synthetic
patients for representative coagulation features, computed as
$|\bar{x}_{\text{synth}} - \bar{x}_{\text{real}}| / |\bar{x}_{\text{real}}|$.}
\label{tab:summary}
\begin{tabular}{lccc}
\toprule
Feature & Visit 1 & Visit 2 & Visit 3 \\
\midrule
Factor II (\%) & 0.018 & 0.017 & 0.007 \\
Factor VIII (\%) & 0.018 & 0.008 & 0.005 \\
AT (\%) & 0.025 & 0.001 & 0.008 \\
Fibrinogen (mg/dL) & 0.005 & 0.009 & $<$0.001 \\
TF Peak (nM) & 0.025 & 0.002 & 0.017 \\
TF ETP (nM$\cdot$min) & 0.013 & 0.008 & 0.023 \\
\bottomrule
\end{tabular}
\end{table}

\clearpage

% ══════════════════════════════════════════════════════════════════════════════
% Figures
% ══════════════════════════════════════════════════════════════════════════════

\begin{figure}[p]
\centering
\includegraphics[width=\textwidth]{validate_bio_correlations.pdf}
\caption{\textbf{Physiological correlations.}
Scatter plots of coagulation factor levels versus thrombin generation parameters
across all three visits. Filled markers are real patients; open markers are
SA-generated synthetic patients. Colors encode visit: blue (V1/baseline),
green (V2/first trimester), orange (V3/third trimester). Real and synthetic
patients occupy the same joint regions across all four pairwise relationships. AT
inversely correlates with TF-initiated peak thrombin (top left), FVIII
positively correlates with peak (top right), and both FII (bottom left) and AT
(bottom right) show the expected relationships with ETP. The visit-dependent
shift in these relationships, reflecting the pregnancy-driven increase in
procoagulant factors, is also preserved in the synthetic cohort.}
\label{fig:bio-correlations}
\end{figure}

\clearpage

\begin{figure}[p]
\centering
\includegraphics[width=\textwidth]{validate_pregnancy_progression.pdf}
\caption{\textbf{Pregnancy-related longitudinal trajectories.}
Mean $\pm$ standard deviation bands for six key coagulation features across
visits (BL = baseline, 1st = first trimester, 3rd = third trimester) for real
(blue) and SA-generated synthetic (orange) patients. The characteristic
pregnancy-driven increases in fibrinogen, Factor~VIII, and vWF are captured, as
are the stable-to-declining patterns in Factor~II and AT.}
\label{fig:pregnancy-progression}
\end{figure}

\clearpage

\begin{figure}[p]
\centering
\includegraphics[width=\textwidth]{validate_cross_visit_corr.pdf}
\caption{\textbf{Cross-visit correlation structure.} Top row: full $216 \times 216$
Pearson correlation matrices for Real ($K{=}23$, left), SA ($N{=}100$, center),
and MVN ($N{=}100$, right) populations. Each matrix is organized as a
$3 \times 3$ grid of $72 \times 72$ blocks (delineated by black lines), where
the diagonal blocks capture within-visit feature correlations (V1--V1, V2--V2,
V3--V3) and the off-diagonal blocks capture cross-visit dependencies (e.g.,
V1--V3 correlations between baseline and third-trimester Factor~VIII).
Stochastic attention preserves both the within-visit and cross-visit block
structure. Bottom row: pairwise residual matrices (element-wise subtraction),
all plotted on the same $\pm 1$ color scale as the top row. The SA$-$Real
residual shows small, unstructured deviations. The MVN$-$Real residual is
smoother, particularly in the off-diagonal blocks, reflecting the Ledoit--Wolf
regularization that systematically shrinks cross-visit correlations toward zero.
The SA$-$MVN residual shows the largest differences in these
off-diagonal blocks, where SA retains longitudinal dependencies that MVN's
rank-deficient covariance estimate cannot represent. A version with amplified
residual color scale ($\pm 0.5$) and Frobenius norms is provided in
\suppfigure{sfig:cross-visit-corr-supp}.}
\label{fig:cross-visit-corr}
\end{figure}

\clearpage

\begin{figure}[p]
\centering
\includegraphics[width=\textwidth]{sim_phase_diagram.pdf}
\caption{\textbf{Population recovery and clinical-cohort subsampling.}
(A)~For each low-rank Gaussian setting, the stochastic-attention memory matrix was constructed
and a sample-covariance Gaussian generator was estimated from the same simulated training
cohort; both generators produced $100$ profiles. Color represented the ratio of their relative Frobenius errors against the known
population covariance (Gaussian error divided by stochastic-attention error), averaged over
latent ranks $3$, $5$, and~$10$. Values above one (blue) indicated lower error for stochastic
attention. Because the target covariance was fixed before the training cohort was drawn,
this measured population recovery rather than reproduction of the training data.
(B)~A stochastic-attention memory matrix was constructed from each of ten random subsets of
the clinical cohort at $n{=}20,15,10,$ and~$8$, and each matrix generated $100$ profiles.
Points and bars represented
the mean and standard deviation of the relative cross-visit correlation error against the
same full $23$-patient reference; the star represented generation using all $23$ patients.}
\label{fig:sim-recovery}
\end{figure}

\clearpage

\begin{figure}[p]
\centering
\includegraphics[width=\textwidth]{validate_conditioned_features_v2.pdf}
\caption{\textbf{Condition-specific feature preservation.} Grouped bar charts
comparing real (dark bars) and SA-generated synthetic (light bars) patient
means ($\pm$ SD error bars) for eight coagulation features, shown separately
for each clinical subgroup: Uncomplicated ($n{=}18$, left, blue), PCOS ($n{=}3$,
center, orange), and Developed~PE ($n{=}5$, right, red). Each bar is the mean of
that assay pooled across all three visits (V1--V3); factor activities (FVIII,
vWF, FIX, FV, $\alpha_2$-AP, FVII, plasminogen) are expressed as \% of normal
pooled reference plasma, and fibrinogen (Fbgn) in mg/dL. Gray percentages
indicate the mean relative error (MRE). Green
annotations show a descriptive bootstrap Mann--Whitney diagnostic: the fraction
of 1,000 bootstrap replicates (synthetic subsampled to $n_{\text{real}}$) in
which $p > 0.05$, meaning that this test detected no difference at the available
sample size. This is not an equivalence test.
Features achieving $\geq$90\% are shown in green; $<$90\% in red. Across all
24 feature--condition pairs, 20 (83\%) had no difference detected in
$\geq$90\% of replicates (median: 98.6\%). The four pairs below this
threshold were FIX and plasminogen in the PCOS and Developed~PE subgroups,
where small subgroup sizes and large inter-patient variability reduced test power.}
\label{fig:conditioned-features}
\end{figure}

\clearpage

\begin{figure}[p]
\centering
\includegraphics[width=\textwidth]{validate_mechanistic_combined_tf_only_v2.pdf}
\caption{\textbf{Mechanistic validation (TF-only).} Top row: BZ2012
ODE-predicted TGA values (vertical axis) versus dataset TGA values (horizontal
axis) for five thrombin generation parameters. The dataset values are the TGA
measurements from each patient's record, present for both real and synthetic
patients. The ODE-predicted values are computed by running each patient's
coagulation factor levels through the 58-species BZ2012 model. The dashed line
indicates where the ODE model would perfectly reproduce the dataset value
($y{=}x$); systematic deviations from this line (e.g., lagtime underprediction)
reflect ODE model bias, not synthetic data failure. Filled markers are real
patients, open markers are SA-generated synthetic patients, colored by visit:
blue (V1/baseline, training set), green (V2/first trimester), orange (V3/third
trimester). Real and synthetic patients occupy the same cloud with the same
pattern of ODE model bias. Bottom row: ODE-predicted/dataset ratio
distributions for real (blue) and synthetic (orange) patients. The
distributions overlap substantially. Each panel is annotated with cloud overlap
(OL), the two-sample Kolmogorov--Smirnov statistic ($D$), and $p$-value. All
five features had $p>0.35$; these tests did not detect a difference between real and
synthetic ratio distributions.}
\label{fig:mechanistic-combined}
\end{figure}

\clearpage

\begin{figure}[p]
\centering
\includegraphics[width=\textwidth]{downstream_utility_scatter_v2.pdf}
\caption{\textbf{Downstream utility: synthetic and real calibrations on held-out real
profiles.} Each panel compares
the BZ2012 ODE predictions for one TGA feature when the model is calibrated on
real V1 patients (horizontal axis) versus synthetic V1 patients (vertical
axis), evaluated on the same held-out real V2 and V3 patients. Points near the
$y{=}x$ line indicate similar predictions from the two calibrations.
Green: V2 (first trimester); orange: V3 (third trimester). The synth-calibrated
model had slightly lower median relative error on all five features (overall
synthetic-to-real error ratio $0.94$). Both calibrations used 11 random restarts; all starts
converged.}
\label{fig:downstream-utility}
\end{figure}
